\section{A certified evaluation on bare terms}
\label{certified}
The same seed, slot arms and allocator admit a second evaluation policy. It acts on bare ground terms, with observations taken between successive tag steps. Between observations it carries one bit and a finite control mode. We describe the boundary map first and then its realization by individual contractions.

Use the following traversals with immutable input. On a word, a pop contracts its first live cell, descending through spent cells at \(LR\). On a slot it searches only the word in a fired one-arm. On a list of slots it tries them in order. On the rail it tries each allocated period in order and skips the frontier. On the complete tree it first tries the header and then the rail. A firing traverses just the rail in the same order. At the first slot with two unfired arms it contracts the arm designated by its bit parameter. An allocation skips the frames and contracts the frontier. At every success the operation rebuilds only the ancestors of the contracted occurrence. On a failed subsearch it continues with the next listed subsearch. The functions are total structural recursions on finite trees. The boundary operation is given by
\[
\begin{array}{ll}
 b\gets h(t).& b=\bot\ \Longrightarrow\ \mathrm{fail}.\\
 t_1\gets\mathrm{pop}(t).&\\
 \text{For }j=0,1,2:&
  \mathrm{fire}_b(t_{j+1})=u\ \Longrightarrow\ \mathrm{return}\ (2+j,u).\\
 &j<2\ \Longrightarrow\ t_{j+2}\gets\mathrm{allocate}(t_{j+1}).
\end{array}
\]
The partial action returns failure when a necessary operation fails. Its total boundary map \(B\) keeps the initial term on failure. A separate verified single-contraction realization has eight modes: boundary, three searches each with one bit, and a terminal fallback. Every invocation begins navigation at the root. The first search fires the chosen arm if present, or allocates and passes to the second search; the second does likewise; the third can fire. Only the bounded stage and the borrowed bit can pass to the next invocation. On an unavailable action the controller tries ordinary preorder, enters the fallback mode on success and reports normality precisely when there is no redex. Lean proves this normality equivalence and one-step legality for every finite input and every mode. For a successful boundary action of reported length \(k\), precisely \(k\) calls take the controller from boundary to boundary at exactly the specified whole ground tree. On the encoded domain every call on a nonempty queue succeeds. An empty queue has already been detected. There is no call to a source interpreter. A source-step counter is used only to state the correctness theorem.

Here \(h(t)\in\{\bot,0,1\}\) is a finite tree-algebra value. On every reachable boundary term it is the first live symbol or \(\bot\). A word observer on \(S\) gives \(\bot\), gives \(b\) on \(J_b w\), and on a spent cell follows the designated left continuation. Combine ordered fields by \(a\triangleright b=a\) if \(a\ne\bot\), and \(b\) otherwise. On a slot look only at the first word inside its fired right arm. Combine these observers over lists and frames, with a final first-choice between header and rail. Constant-depth shape tests give a total interpretation on arbitrary trees. The actual bottom-up state has five Boolean and fifteen ternary fields, including the promised head of a word, a slot, a row, a rail and the whole tree. The other fields store bounded-depth pattern matches on the two children. The finite transition table, expressed as constructor equations, is also given in the symbolic supplement. Lean checks an injection of this state into twenty finite coordinates. The detector is the single fixed regular language \(H(t)=[h(t)=\bot]\).

\begin{theorem}\label{certmain}
Let \(B\) be the bare boundary function and \(V\) the structural live-word projection. For every nonempty finite binary cyclic program, every initial word and every \(n\), \(V(B^nE(P,w))\) is the source queue after \(n\) absorbing source steps. On every finite tree, every successful action is a chain of exactly its reported \(k\) contextual \(S\) contractions, where \(2\le k\le4\). For the explicit fixed tape-to-Cook compiler \(C\), fixed program \(P_*\), and every finite deterministic binary tape-machine instance \(M\),
\[
             M\downarrow\quad\Longleftrightarrow\quad
              \exists n\in\mathbb N:\ H\!\left(B^n E(P_*,C(M))\right)=1.
\]
All three statements are checked in Lean.
\end{theorem}
\begin{proof}
Represent one visible word by its finite sequence of spent bits and opaque saved tails followed by a live word. A positive spent slot surrounds exactly one such editable segment, with the unedited right duplicate opaque; a zero slot has no editable segment. A completed frame is the left context wrapping a packed list of these spent slots around its terminal \(S\) and a pristine allocation stamp. A well-formed outer list is one editable header followed by completed frames. In an open period, a second list of spent slots ends at the still-pristine suffix of the program. These are the inductive cases \texttt{SlotOK}, \texttt{FrameOK}, \texttt{OuterOK} and \texttt{Valid}.

For a word with spent prefix \(h\), the pop search either fails on an empty live suffix or gives the term with its first bit newly spent. Induct on \(|h|\), rebuilding the left continuation at each return. In a positive slot, rebuild the box, the open arm and the slot. In a zero slot take the next branch. Induction on the slot list and then on the outer frames shows that popping the bare term gives the same tree as erasing the annotated pop. The same inductions show the observable head agrees with the first live letter of the physical collection. Popping preserves its grammar: it changes only the indicated segment, and the copied history stays opaque. This also proves actual word readback including its chronological spent prefix.

A fire search skips all completed slots. If an open period is present, it therefore reaches its first pristine slot. The two arms give precisely the two terms of \eqref{commit}, and the remaining slots keep their order. At a frontier, firing fails, and allocation obeys \eqref{reset} or \eqref{alloc}. For the zero case the next frontier is \(X_{j+1}X_{j+1}\); the following allocation succeeds. For the positive case allocation succeeds immediately. The newly exposed period is well-formed. Firing either extends the open list by one spent slot or moves the completed period to the outer list. Thus every case preserves the grammar, and the unannotated action and annotated advancement return equal contraction counts and identical whole ground trees. Induction on \(n\) proves their iterates have equal erasures, including absorption at an empty word.

Independently, induction on each raw traversal shows that any success is a root \(S\) contraction in a one-hole context, even for inputs outside the grammar. The three possible successful branches have lengths \(2,3,4\). The compiler \(C\) first translates the transition table and initial tape into a primitive three-counter program and its initial registers. It next compiles that job to a well-formed restricted deletion-two tag job. It encodes the job for the Rogozhin interpreter and the fixed period-912 Cook system. The literal functions are \texttt{compileThreeCounterJob}, \texttt{compileT2Job} and \texttt{encodeBits} in Filin's \texttt{DeterministicTapeCook} module. The checked equivalences at these three interfaces compose to identify source halting with eventual cyclic dataword emptiness. These compilers terminate without using the source's halting time. Composing this theorem with the preceding identities proves the displayed equivalence. The formal theorem depends only on \texttt{propext} and \texttt{Quot.sound}.
\end{proof}

\begin{theorem}\label{microhalt}
Let \(\Delta\) be one call of the eight-mode controller, extended by the identity on normal terms. Put \(J(m,t)=[m=\mathsf{boundary}]H(t)\). For every deterministic tape-machine instance \(M\),
\[
 M\downarrow\quad\Longleftrightarrow\quad
 \exists k:\ J\!\left(\Delta^k(\mathsf{boundary},E(P_*,C(M)))\right)=1 .
\]
The guard may be tested after every individual contraction. It is a fixed regular test with one finite mode observation.
\end{theorem}
\begin{proof}
In a successful block of length \(l=2,3,4\), the interior modes are respectively the first \(l-1\) search stages with the same borrowed bit. Its endpoint is the boundary mode with exactly the tree returned by \(B\). Thus no interior stage accepts. Induction on the macro index realizes any first successful macro observation at some finite primitive time. Conversely, at any generated boundary with \(H=0\), the physical queue is nonempty and the next block exists. If a primitive-time acceptance occurs \(k\) steps later, either \(k=0\) and the initial boundary accepts, or \(k\) is at least the next block length \(l\). In the latter case shift to the resulting boundary with \(k-l<k\). Strong induction gives an accepting macro observation. An already accepting boundary needs no successor condition. Theorem~\ref{certmain} now gives the result. The complete argument is the declaration \texttt{UniformS.Micro.universal\_single\_contractions}.
\end{proof}

\paragraph{\textit{Finite output.}}
For \(d\) source machines, encode the \(d\) seeds in a fixed list of \(S\)-pairs. One round advances each branch once, rebuilding its unaltered siblings. Context closure gives at most \(4d\) native contractions per round. The \(i\)th detector is the preceding regular language at the \(i\)th list field. By induction its first success exists precisely when machine \(i\) halts. For two recognizers with disjoint halting domains, accept when either detector succeeds and return its index. The output is the literal identity of the accepting port. Every partial computable Boolean function admits this standard presentation: run its computing machine and accept precisely when its result is the designated bit, with a permanent loop on the opposite result. The two recognizers are finite programs, fixed independently of their input. The two-port statement is also formalized in Lean. This result differs from the single numerical menu of Theorem~\ref{output}, which returns an unbounded unary count from one physical queue.

\paragraph{\textit{Encoding and control costs.}}
Set \(m(0)=3,\ m(1)=4\), extending \(m\) additively to words. Unfolding \eqref{seed} and using structural induction gives the following kernel-checked exact leaf count:
\[
 L(E(P,w))=22+48p+4\sum_{i<p}m(v_i)+m(w).
\]
The bounded protocol can start each constituent search at the root. It passes only a bit and the finite stage between these searches; nothing is retained between macro calls. The bottom-up detector uses five bits and fifteen ternary values per summary. Direct evaluation with a recursion stack needs \(O(h)\) such summaries for a height-\(h\) occurrence tree. Each search is a linear-time scan with a bounded number of visits to each relevant occurrence. Parent links support root-restarted navigation as in Proposition~\ref{walker}. The direct Lean evaluator rebuilds persistent context, so its call stack and constructed term also consume space. These costs differ from the bounded number of contractions. On replacement systems of Theorem~\ref{cutoff} the same regular halt observations, even with feedback, yield a bounded abstract orbit. The concrete implementation admits no such replacement for its universal halting contract.
